% ======================================================================
% Эксперименты 3.5.3–3.5.5. Ответственный: C — Дучин Кирилл.
%
% Что сдать (main.tex: теория 1518–1536, пункты 1547–1557, ожидаемый результат
% 1568–1575, пример 1626–1699; требования к отчёту 1875–1879):
%   3.5.3  Симулятор системы методом Эйлера; найти шаг Delta t потери
%          устойчивости; сравнить с условием |1+Delta t*nu|<1 для собственных
%          чисел из п. 3.3; траектории q(t), c(t) при указанных начальных условиях.
%   3.5.4  Численно найти порог флаппинга для u0=0.6; 0.75; 0.9 и сравнить с k*;
%          описать критерий обнаружения; проверить устойчивость наблюдения
%          при уменьшении временного шага.
%   3.5.5  Градиентный спуск для Phi(c1,c2) из п. 3.4 со стартовой точки (10,10)
%          при постоянном и убывающем шаге; число итераций, точка, значение
%          функции; траектория на фоне линий уровня; сравнение с аналитикой.
%          ОБРАЗЕЦ ОФОРМЛЕНИЯ ЭТОГО ЭКСПЕРИМЕНТА — report_template.tex, 203–336.
%
% Параметры варианта 3 (приложение Б): lambda=600, mu=80, mu2=120, a=1, b=1.
% NB (main.tex, строки 1692–1699): перед запуском проверить допустимость (10,10):
% mu_i*c_i > lambda; если недопустима — использовать (R1+1, R2+1) и указать замену.
% Зафиксировать начальные условия симуляции, длительность и критерий обнаружения
% колебаний. Численный градиент — переиспользовать функции разностей из п. 3.5.1
% (блок A в solution.py), не дублировать.
%
% Запуск: python experiments.py 3.5.3 / 3.5.4 / 3.5.5.
% Рисунки: figures/3-5-3-trajectories.pdf, figures/3-5-5-paths.pdf.
% ======================================================================

\taskhead{Эксперимент 3.5.3. Моделирование методом Эйлера и граница устойчивости}

\taskhead{Постановка эксперимента}
\todo{C: описать симулятор (уравнения, начальные условия, длительность, перебираемые $\Delta t$); указать, как фиксировалась потеря устойчивости.}

\taskhead{Программный результат и сравнение}
\todo{C: листинг вызова своих функций из solution.py и запуск python experiments.py 3.5.3.}
\taskfigure[0.85]{figures/3-5-3-trajectories.pdf}{Траектории $q(t)$ и $c(t)$ при выбранных начальных условиях и настройках}{fig:3-5-3}
\begin{table}[htbp]
\centering\small
\caption{Граница устойчивости по временному шагу: численный порог и линейный критерий}
\label{tab:3-5-3}
\begin{tabular}{|r|r|r|r|r|}\hline
$u_0$ & $k$ & $\Delta t$ (числ. порог) & $2/|\nu|$ (критерий) & Расхождение\\\hline
— & — & — & — & —\\
\hline\end{tabular}
\end{table}
% Расхождение численного порога и линейного критерия — не «ошибка»:
% критерий относится к линейному приближению вблизи равновесия
% (main.tex, строки 1523–1527). Пояснить причину расхождения.
\todo{C: пояснить график и таблицу; сравнить с условием $|1+\Delta t\,\nu|<1$ для $\nu$ из п. 3.3.}

\taskhead{Эксперимент 3.5.4. Численный порог флаппинга}

\taskhead{Постановка эксперимента}
\todo{C: описать критерий обнаружения колебаний (начальные условия, длительность, допуск, перебор $k$); одно пересечение уровня $c^*$ не доказывает колебаний — анализировать повторные смены направления отклонения (main.tex, строки 1652–1661).}

\taskhead{Программный результат и сравнение}
\todo{C: листинг вызовов и запуск python experiments.py 3.5.4.}
\begin{table}[htbp]
\centering\small
\caption{Пороги колебательного режима: аналитический $k_*$ и численный порог}
\label{tab:3-5-4}
\begin{tabular}{|r|r|r|r|}\hline
$u_0$ & $k_*$ (п. 3.3) & Численный порог & Расхождение\\\hline
0.6 & — & — & —\\
0.75 & — & — & —\\
0.9 & — & — & —\\
\hline\end{tabular}
\end{table}
\todo{C: пояснить расхождения (зависимость от начальных условий, шага, длительности, допуска); проверить устойчивость наблюдения при уменьшении $\Delta t$.}

\taskhead{Эксперимент 3.5.5. Градиентный спуск для функции затрат}

\taskhead{Постановка эксперимента}
\todo{C: проверить допустимость $(10,10)$ для варианта 3 (при нарушении — $(R_1+1,R_2+1)$ и указать замену); описать численный градиент (схема, шаг $h_d$), два правила шага $\eta_j$ (постоянный и убывающий — выбрать и записать), критерий остановки, предел итераций, проверку допустимости на каждом шаге (образец: \texttt{report\_template.tex}, строки 237–261).}

\taskhead{Программный результат и сравнение}
\todo{C: листинг вызовов и запуск python experiments.py 3.5.5.}
\begin{table}[htbp]
\centering\small
\caption{Результаты градиентного спуска}
\label{tab:3-5-5}
\begin{tabular}{|l|r|r|r|r|r|}\hline
Шаг & Итерации & $\widehat c_1$ & $\widehat c_2$ & $\Phi(\widehat c)$ & $\Delta_c$\\\hline
Постоянный & — & — & — & — & —\\
Убывающий & — & — & — & — & —\\
\hline\end{tabular}
\end{table}
\taskfigure[0.85]{figures/3-5-5-paths.pdf}{Траектории спуска на фоне линий уровня затрат; отмечены начальная точка и аналитический минимум}{fig:3-5-5}
\todo{C: сравнить с аналитическим минимумом из п. 3.4 ($\Delta_c$, $\delta_c$ — определения в \texttt{report\_template.tex}, строки 274–279); ошибки затрат; допустимость всех точек.}

\taskhead{Вывод}
\todo{C: сравнение правил шага (осторожно: вывод относится только к выбранным правилам и начальной точке); причины расхождений.}
